\documentclass[12pt]{book}
\usepackage[utf8]{inputenc}
\usepackage[T1]{fontenc}
\usepackage{amsmath,amssymb,amsfonts}
\usepackage{mathrsfs}
\usepackage{amsthm}

% Calligraphic for categories and sheaves
\newcommand{\cat}[1]{\mathcal{#1}}
\newcommand{\sheaf}[1]{\mathcal{#1}}

% Standard math operators
\DeclareMathOperator{\Hom}{Hom}
\DeclareMathOperator{\Ext}{Ext}
\DeclareMathOperator{\Tor}{Tor}
\DeclareMathOperator{\id}{id}
\DeclareMathOperator{\varinjlim}{\varinjlim}

\begin{document}

\setcounter{page}{275}
\section{5. Dualizing functors on a local noetherian ring.}

In this section we recall without proof the definition and some properties of dualizing functors on a local noetherian ring \(A\), with maximal ideal \(\mathfrak{m}\).
For proofs, see [LC,84].

\begin{proposition}
Let \(I\) be an injective hull of the residue field \(k\) of \(A\), and denote by \(T\) the functor \(\Hom(\cdot,I)\).
Then for every \(A\)-module \(M\) of finite length, the natural map
\[M \to T T(M)\]
is an isomorphism.
\end{proposition}

\textbf{Definition.}
A contravariant additive functor \(T\) from the category \(c_{f}\) of \(A\)-modules of finite length into itself is called a dualizing functor for \(A\) if there exists an injective hull \(I\) of \(k\) and an isomorphism of functors \(T \cong \Hom(\cdot, I)\).

\begin{proposition}
A contravariant additive functor \(T\) from \(c_{f}\) into itself is a dualizing functor if and only if it is exact, and \(T(k) \cong k\).
In that case one can take \(I=\varinjlim_{n} T(A/\mathfrak{m}^{n})\) and then there is a canonical isomorphism \(T \cong \Hom(\cdot, I)\).
\end{proposition}

\begin{proposition}
If \(I\) is an injective hull of \(k\), then the functor \(\Hom(-,I)\) gives an anti-isomorphism of the category (dcc) of \(A\)-modules with descending chain condition (we call them \(A\)-modules of co-finite type) and the category (acc) of modules of finite type over the completion \(\hat{A}\) of \(A\).
\end{proposition}

\end{document}